\chapter{Controls and Indicators}
\label{ch:controls}

You drive the machine with three things: the keypad, the mouse, and the
display. The keypad edits recipes and runs utilities; the mouse runs bonds; the
display tells you what is happening.

\figureplaceholder{Photograph of the control panel with the left navigation
keys, the right machine buttons and the four indicator lamps labelled.}

\section{The mouse}

The mouse has two buttons and is how you actually bond. It does not move a
cursor.

\begin{center}\small
\begin{tabular}{@{}L{26mm}L{95mm}@{}}
\toprule
\textbf{Button} & \textbf{What it does} \\
\midrule
Right & Starts a bond cycle, and triggers each phase within it. In Manual mode,
        also drives the head \emph{up}. \\
Left  & Drives the head \emph{down} in Manual mode. Does nothing in the other
        modes. \\
\bottomrule
\end{tabular}
\end{center}

\begin{wbnote}
The right button is used both for pressing and for releasing. In the automatic
modes the press starts the phase and the \emph{release} begins the descent ---
so the machine waits for you to let go before it commits to moving. Hold the
button while you position yourself, release when you are ready.
\end{wbnote}

\begin{wbnote}
The right button is also what you hold during a force calibration run
(Chapter~\ref{ch:calibration}), not the left.
\end{wbnote}

\section{The keypad}

% What every key does on every page.
%
% DERIVED FROM FIRMWARE. Refresh against:
%   Firmware/Core/Src/menu.cpp                  (handleUp/Down/Left/Right/Plus/
%                                                Minus/Save/Load/Enter/Add/
%                                                EscapeDelete, lines 342-531)
%   Firmware/Core/Src/user_interface_module.cpp (hotkey targets, machine buttons)
%   Firmware/Core/Inc/configuration.h           (KEYPAD_BTN_REPEAT_*)
%
% A blank cell means the key does nothing on that page.

\newcommand{\keyrow}[7]{\keycap{#1} & #2 & #3 & #4 & #5 & #6 & #7 \\}

\begin{center}\footnotesize
\begin{tabular}{@{}L{18mm}L{20mm}L{20mm}L{18mm}L{20mm}L{16mm}L{16mm}@{}}
\toprule
\textbf{Key} & \textbf{Parameter Edit} & \textbf{Configuration Select} &
\textbf{Settings} & \textbf{Name entry} & \textbf{Delete confirm} &
\textbf{Force Setup} \\
\midrule

\keyrow{$\uparrow$}{pointer up}{pointer up}{pointer up}{next character}{select \menuitem{Y}}{value up}
\keyrow{$\downarrow$}{pointer down}{pointer down}{pointer down}{previous character}{select \menuitem{N}}{value down}
\keyrow{$\leftarrow$}{previous screen}{previous mode}{}{cursor left}{select \menuitem{Y}}{}
\keyrow{$\rightarrow$}{next screen}{next mode}{}{cursor right}{select \menuitem{N}}{}
\addlinespace
\keyrow{$+$}{increase value}{}{increase value}{next character}{}{value up}
\keyrow{$-$}{decrease value}{}{decrease value}{previous character}{}{value down}
\addlinespace
\keyrow{ENTER}{}{load pointed}{toggle / run row}{confirm}{apply answer}{save}
\keyrow{LOAD}{open selector}{load pointed}{}{}{}{}
\keyrow{SAVE}{save configuration}{open Settings}{save settings}{}{}{save}
\keyrow{ADD}{}{new configuration}{}{}{}{}
\keyrow{ESC/DEL}{open selector}{delete pointed}{back}{cancel}{cancel}{cancel}
\addlinespace
\keyrow{TAIL $\pm$}{edit \menuitem{Tail}}{}{}{}{}{}
\keyrow{LOOP $\pm$}{edit \menuitem{Loop Height}}{}{}{}{}{}
\keyrow{SEARCH $\pm$}{edit \menuitem{Search 1}}{}{}{}{}{}
\keyrow{STEP $\pm$}{edit \menuitem{Stepback}}{}{}{}{}{}

\bottomrule
\end{tabular}
\end{center}

\begin{wbnote}
On the Settings page \keycap{ENTER} acts on the pointed row: it toggles
\menuitem{Spot On}, and starts tachometer calibration on
\menuitem{Start Tach.\ Cal.} On every other row it does nothing.

In Table Tear mode \keycap{TAIL $\pm$} edits \menuitem{Table Tail} rather than
\menuitem{Tail}. If a hotkey targets a parameter the active mode does not use,
the machine shows \msg{Not used by protocol}.
\end{wbnote}

% ------------------------------------------------------------ machine keys --

\subsection*{Machine buttons}

These act on the machine itself and work from any page.

\begin{center}\small
\begin{tabular}{@{}L{26mm}L{95mm}@{}}
\toprule
\textbf{Button} & \textbf{Action} \\
\midrule
\keycap{SETUP}      & Starts a force calibration run. Refused with
                      \msg{BUSY! TRY LATER} if the machine is running. \\
\keycap{TEST}       & Starts an ultrasonic test and reports frequency, Q and
                      power. Refused while busy. \\
\keycap{CLAMP OPEN} & Opens or closes the wire clamp. Refused while busy. \\
\keycap{LIGHT}      & Turns the area light on or off. Works at any time. \\
\keycap{MANUAL}     & Switches the active configuration to Manual mode. \\
\keycap{RESET}      & Restarts the machine immediately. See the warning below. \\
\bottomrule
\end{tabular}
\end{center}

\begin{wbwarning}
\keycap{RESET} is an immediate hardware restart, not a soft stop. It takes
effect at once, from any state, including mid-weld. Motion stops wherever it
is and the machine re-homes on restart. Use it to clear a locked machine, not
to end a cycle.
\end{wbwarning}

\begin{wbnote}
The panel also carries a \keycap{HIGH RESET} button. It is not connected to any
function in this firmware version and does nothing when pressed.
\end{wbnote}

% ----------------------------------------------------------- key repeating --

\subsection*{Holding a key down}

The value keys --- $\uparrow$, $\downarrow$, $+$, $-$ and the four hotkey pairs
--- repeat while held, and accelerate the longer you hold them:

\begin{center}\small
\begin{tabular}{@{}L{40mm}L{40mm}L{40mm}@{}}
\toprule
\textbf{Held for} & \textbf{Repeat rate} & \textbf{Step size} \\
\midrule
under 0.5\,s      & ---                  & one step per press \\
0.5\,s to 1.5\,s  & 4 per second         & 1 step \\
1.5\,s to 4.5\,s  & 4 per second         & 10 steps \\
over 4.5\,s       & 4 per second         & 100 steps \\
\bottomrule
\end{tabular}
\end{center}

So a parameter with a step of 0.01\,mm moves at 0.04\,mm/s at first, 0.4\,mm/s
after a second and a half, and 4\,mm/s after four and a half. Release and press
again to return to single steps.

\begin{wbnote}
$\leftarrow$ and $\rightarrow$ do \emph{not} repeat except while entering a
name. Holding them will not page through parameter screens or bonding modes ---
press once per step.
\end{wbnote}

\derivedfrom{\texttt{menu.cpp}, \texttt{user\_interface\_module.cpp},
             \texttt{configuration.h}}


\section{The indicator lamps}

Four lamps report machine state. All four are off at power-up.

\begin{center}\small
\begin{tabular}{@{}L{26mm}L{95mm}@{}}
\toprule
\textbf{Lamp} & \textbf{Lit when} \\
\midrule
\keycap{TEST}       & An ultrasonic test is running. \\
\keycap{SETUP}      & A force calibration run is in progress. \\
\keycap{CLAMP OPEN} & The wire clamp is open. \\
\keycap{MANUAL}     & The loaded configuration is a Manual-mode one. \\
\bottomrule
\end{tabular}
\end{center}

\begin{wbnote}
The \keycap{TEST} and \keycap{SETUP} lamps follow what the machine is actually
doing, not what you asked for. If a run fails, its lamp goes out --- it cannot
be left stranded on.
\end{wbnote}

\section{Lighting}

The \keycap{LIGHT} button switches the area light on and off. Its brightness,
and the spotlight's, are set on the Settings page
(Chapter~\ref{ch:settings}). The spotlight has no panel button; it is switched
on the Settings page too.
